\section{What fog and rain do to the perception: the measurement}\label{sec:perception}

\begin{plainwords}
A real car would see through cameras. Fog and rain make cameras worse. Before testing whether weather hurts a driving model, we measured how much the weather filter of the project really damages the two perception steps that feed the model: the object detector that finds vehicles, and the stereo-depth algorithm that estimates how far they are. The result is that the damage is large, and that it consists mostly of \emph{missed vehicles}.
\end{plainwords}
\begin{cautionbox}[Compared with GitHub]
The functions measured here are the teammates' modules committed on GitHub (\pth{YOLO+weather module/}, \pth{bev_state/}); they were only called, never modified. The GitHub README describes the weather filter as ``preliminary tooling'' that ``is not evidence of physical fog/rain robustness''. GitHub contains \emph{no measurement} of the filter's effect on detection or depth; the measurement in this section is local only (Section~\ref{sec:github}).
\end{cautionbox}


\subsection{Why it was measured}
The teammates' weather module corrupts the camera images (Koschmieder-style fog with per-pixel depth, alpha-composited rain streaks), and the stereo depth that places a detected vehicle in the grid is computed from the corrupted image pair.
An earlier policy-side proxy removed vehicles from the grid at random with a probability taken from an older measurement of the detector (``detector dropout''). Every policy was completely flat under that proxy, so no claim about adaptation could be tested.
The project director's position was that weather also degrades the \emph{depth}; this had to be measured before anything was trained on it. The measurement (\pth{bev_rl/adaptive/measure_depth_error.py}) only \emph{calls} the teammates' functions and compares with the MetaDrive ground truth; it modifies nothing in the teammates' modules.
Hypotheses and decision rules were written down \emph{before} the measurement (\pth{report/data/preregistration_depth_degradation.md}).

\subsection{Protocol}
Seeds 7000--7039 (40 scenes, traffic density 0.3), stationary ego car, traffic vehicles placed 5--19\,m ahead, each scene perceived clean and with two random weather draws per condition. Conditions: clear, fog 0.6, 0.8, 1.0 and rain 0.3, 0.6, 1.0
(the number is the intensity parameter of the teammates' augmentation). For every vehicle in view the following were recorded: whether YOLO detected it, whether the stereo depth at the vehicle's pixel was valid, and the range error of the box depth.

\subsection{Result}
\explainedfigure{fig_measured_perception_degradation}{0.97}{Measured effect of simulated fog and rain on the camera perception}{Left: the share of vehicles in view that the YOLO object detector finds, per weather condition (vertical lines are 95\% intervals). Right: the share of vehicle pixels at which the stereo-depth algorithm returns a valid distance. Weather conditions are the teammates' image augmentation at the intensity shown (0 = none, 1 = strongest). 40 scenes; 72 vehicle observations in clear weather and 144 in each weather condition.}{In clear weather the detector finds 72\% of the vehicles. In fog of intensity 1.0 it finds 31\%, in rain of intensity 0.6 it finds 15\% and in rain of intensity 1.0 only 1.4\%. Stereo depth stays valid for 88\% of pixels in clear weather and 49\% in the strongest fog, but remains at 67\% to 83\% in rain.}{The perception measurement confirms that the simulated weather is a real degradation of the camera pipeline, mostly through \emph{missed vehicles}. For vehicles that are detected, the distance estimate is not worse. This tells us what a driving model would have to cope with, and explains why the driving results in the next figures are unaffected.}{fig:perception}


\begin{itemize}
\item \textbf{Detection collapses with weather} (hypothesis H0 in the pre-registration is confirmed): 0.72 in clear weather, 0.72 in fog 0.6, 0.59 in fog 0.8, 0.31 in fog 1.0, 0.57 in rain 0.3, 0.15 in rain 0.6 and 0.014 in rain 1.0.
\item \textbf{Stereo depth validity drops in fog} (0.88 clear, 0.63 in fog 0.8, 0.49 in fog 1.0) but much less in rain (0.81 at rain 0.6, 0.67 at rain 1.0).
\item \textbf{For vehicles that are detected, the range estimate does not get worse}: the box depth was valid for 100\% of detections in every condition and the median absolute depth error is \emph{lower} in fog and rain (0.98\,m clear, 0.36\,m fog 1.0, 0.67\,m rain 1.0). The share of detections with a relative error above 20\% does not rise (0.096 clear; 0.091 at fog 1.0; 0 in rain).
  The likely explanation is a selection effect (hypothesis, not tested): only close, large vehicles survive detection in heavy weather.
\item Consequently the degradation that reaches a policy is almost entirely \emph{missed detections}, and the ``wrong range'' branch of the calibrated environment is rarely triggered. This is the reason why even the calibrated degradation does not hurt the policy in sparse traffic (Section~\ref{sec:weather}).
\end{itemize}

\subsection{How detector retention was measured at different intensities, and how the driving model receives it}
This is the measurement that actually feeds the driving models in every weather test of this report. \pth{bev_rl/adaptive/measure_paired.py} placed up to three vehicles around a stationary car in 20 scenes (seeds 1000--1019), perceived each scene clean and with two random weather draws at each of six fog and six rain intensities, and counted the detector's boxes. The result is a table of retention values (\pth{bev_rl/adaptive/results/retention_table.json}). During a weather test the driving model's grid is built from the simulator's vehicle list, and each vehicle is dropped with probability one minus the measured retention at the tested intensity (linear interpolation between measured points, retention capped at 1).
\explainedfigure{fig_detector_retention_by_weather_intensity}{0.95}{Detector retention measured at different fog and rain intensities, and the probability that the driving model keeps a vehicle}{In 20 scenes (seeds 1000 to 1019, traffic density 0.3, stationary car, up to three vehicles placed around it) the same camera frames were passed through the teammates' vehicle detector twice: once clean and once with the weather filter at each intensity (two random weather draws per scene). Retention is the number of vehicles found with weather divided by the number found without weather. Grey dots are single scenes, the blue line is the mean over scenes with its 95\% interval, the orange dashed line is the measured retention capped at 1, which is the probability that the weather proxy keeps a vehicle in the driving model's grid. Dotted vertical lines mark the intensities used in the 200-episode weather tests.}{Fog does almost nothing up to intensity 0.6, then removes vehicles quickly: about 71\% are kept at 0.8 and 27\% at 1.0. Rain acts earlier and harder: about half of the vehicles are kept at 0.45, 16\% at 0.6 and 4\% at 1.0. Values slightly above 1 in light fog mean that the detector found a few more boxes than in the clean frame; they are treated as no loss.}{The driving model never sees the weather and never sees camera images in these tests. It receives the weather only indirectly: each vehicle is removed from its bird's-eye-view grid with the probability read off this curve. The conclusions about weather therefore hold for \emph{loss of detected vehicles at these measured rates}, from 20 scenes per point; they do not cover other perception failures.}{fig:detectorretention}


\subsection{Consistency with the earlier retention table}
The detector-dropout proxy used a retention table measured earlier (\pth{bev_rl/adaptive/results/retention_table.json}): fog 1.0, rain 0.6 and rain 1.0 keep 0.27, 0.16 and 0.04 of the vehicles. The corrected measurement above gives 0.31, 0.15 and 0.014, i.e.\ the same order of magnitude;
the table's bias suspected after the stale-camera error (below) therefore seems small for these three conditions, but it has not been re-fitted.

\subsection{An invalid first measurement (kept for the record)}
The first run of the measurement (20:31--20:52 on 10 November) was \textbf{invalid}. It called a helper that mounts the camera rig \emph{once per environment object} after every reset, so scenes 2--40 had stale cameras: stereo depth was valid for only 2 of 73 clean vehicle-frames and was identical in all weather conditions.
The teammates' own benchmark script remounts the rig after every reset; the corrected run does the same. The invalid outputs were renamed \code{*.SUSPECT_*} and are kept; the pipeline stage that had already used the invalid model was discarded and rerun. The same helper had been used by the earlier measurement that produced the retention table, which is why that table carries a caveat.
\emph{Lesson:} a perception measurement must be validated by an ``is the effect identical across conditions'' check before use.

\subsection{Cost of the full perception chain}
Benchmark of the teammates' chain (30 steps, seed 1000, no traffic): full resolution 1200$\times$900 with YOLO costs 1.7\,s (clear), 2.1\,s (fog 0.6), 2.5\,s (rain 0.6) per step; rendering the cameras at 600$\times$450 with a cached fog transmission, a cropped morphology region and YOLO disabled when there is no traffic
brings this to 0.20, 0.39 and 0.52\,s. Road-surface accuracy of the half-resolution variant was comparable (0.946 versus 0.945). Even so, a 100\,000-step reinforcement-learning run through this chain would take 50 hours or more, which is why the camera-free surrogate is used (Section~\ref{sec:surrogate}).
